% Ticker, ISIN und Datum sind Bezeichner, keine Kennzahlen:
% check-literals: skip
\begin{titlepage}
\centering
\vspace*{3cm}
{\Huge\bfseries Neste Oyj\par}
\vspace{0.8cm}
{\Large Anlageanalyse aus Prim\"arquellen\\[0.2cm] mit zeitlicher Reichweite\par}
\vspace{0.4cm}
{\large Nasdaq Helsinki: NESTE \textbar{} ISIN FI0009013296\par}
\vspace{2.5cm}
{\large Kurs \je{\Kurs}{EUR} am \KursDatum\par}
\vspace{0.3cm}
{\large B\"orsenwert \Mio{\Boersenwert}{EUR}\par}
\vspace{3cm}
{\large Stand: 27.~September 2026\par}
\vspace{1cm}
\begin{minipage}{0.86\linewidth}\small
\textbf{Die Frage, die dieser Bericht beantwortet:} Soll ein Anleger die Neste-Aktie zum
heutigen Kurs kaufen, halten oder meiden, und wie h\"angt die Antwort vom Anlagehorizont
ab: zw\"olf Monate, drei Jahre (bis nach der Inbetriebnahme in Rotterdam) oder f\"unf bis
f\"unfzehn Jahre?

\medskip
\textbf{Farben.} Schwarz sind belegte Tatsachen und Rechnungen. \bk{Blau ist die
Einsch\"atzung des Verfassers.} \vd{Rot ist das Votum.} Jede Zahl im Flie{\ss}text stammt
aus der Zahlenschicht \texttt{data/}; \GeprueftAnzahl{} von \GeprueftGesamt{} gespeicherten
Prim\"arwerten sind gegen die gedruckte Seite gepr\"uft, die sie nennen.
\end{minipage}
\end{titlepage}
